\chapter*{Conclusioni}
\label{chap:conclusioni}
\phantomsection
\addcontentsline{toc}{chapter}{Conclusioni}
\markboth{Conclusioni}{Conclusioni}

% Riscrittura del 2026-11-19 sulla nuova struttura dei capitoli 5 e 6 (docs/STRUTTURA_NARRATIVA_E_RISULTATI_CAPITOLI_5_6_7.md,
% sezione 6). Testo unico senza suddivisioni, non numerato come l'introduzione (richiesta dell'utente). Numeri ripresi dai
% capitoli 3-6, stesse fonti. DA AGGIORNARE all'esito dello scan delle relazioni locali prima del pooling (Legion
% 2123991-2123995): il paragrafo sulle prospettive del modello lo cita come verifica in corso.

Questa tesi è partita da una domanda semplice da porre e meno semplice da affrontare, cioè come convenga spendere il budget di uno Spiking Transformer con circa mezzo milione di parametri perché colga le relazioni temporali che contano nei dati a eventi, senza perdere di vista quanto costi ricordare. La risposta non è venuta da un singolo modulo, ma da una sequenza di scelte, ciascuna misurata cambiando un solo elemento alla volta, e conviene ripercorrerla nell'ordine in cui si è costruita.

La prima scelta ha riguardato la baseline, una versione compatta di QKFormer in cui quasi tutto il costo si concentrava nei due patch embedding, mentre la memoria temporale era affidata ai soli potenziali di membrana. Spostare risorse dal frontend a una memoria esplicita, con un frontend piramidale, un mixer causale a più ritardi e una convoluzione depthwise al posto dell'attenzione del primo stadio, ha portato la Macro-F1 di DVS-Lip da 44 a 55 punti con appena 320 parametri in più e con il $42\%$ di stato persistente in meno. Il vantaggio si è ripetuto sui quattro dataset, e la baseline non lo raggiunge nemmeno con un budget quattro volte più grande. In un modello compatto conta dove la capacità viene spesa, prima ancora di quanta se ne aggiunga, ed è questa la prima indicazione del lavoro, da cui in qualche misura discendono tutte le altre.

Il modello così ottenuto contava il tempo in passi, e la seconda scelta ha riguardato proprio il significato di un passo. Sostituire la finestra fissa con una rappresentazione a lunghezza variabile, in cui ogni passo ha la stessa durata fisica in tutto il dataset, ha fatto guadagnare oltre nove punti su THU-EACT-50-CHL e circa otto su DVS128 Gesture, un effetto paragonabile a quello di una modifica dell'architettura. Un controllo con un numero fisso di passi ha mostrato che eliminare le code vuote della finestra fissa non basta, e che dove le durate delle registrazioni variano molto conta la costanza fisica del passo. Il clock dell'input, spesso trattato come un dettaglio di preparazione dei dati, si è rivelato una variabile di progetto a tutti gli effetti, e ha reso possibile ciò che è venuto dopo, perché con un passo di durata nota i ritardi della rete diventano distanze precise nel passato.

Su quelle distanze poggia il contributo centrale, nato dall'osservazione che il mixer a ritardi del modello di riferimento richiamava il passato e lo sommava al presente, senza mai confrontarli. L'IRM antisimmetrico riassume il passato in bande a distanze fisiche fisse, tiene conto di quanta storia esista davvero e confronta ciascun riassunto con il presente attraverso un termine bilineare a basso rango, che cambia segno quando i due istanti si scambiano, sulla falsariga del correlatore di Hassenstein e Reichardt. Su THU-EACT-50-CHL, con mezzo milione di parametri, supera il modello di riferimento di quasi otto punti di Macro-F1 in validazione e di quasi undici sul test ufficiale. Né un sito vuoto né un percettrone con gli stessi parametri riproducono il guadagno, e le stesse bande senza confronto restano al livello del riferimento. Il merito è del confronto, che nella forma antisimmetrica rende più che in quella simmetrica, applicato a un passato collocato a distanze fisiche precise. Tracce esponenziali, un orologio appreso e un passato più remoto, provati al posto delle bande, perdono tutti terreno. Gli interventi sul modello addestrato raccontano la stessa storia dall'interno, perché la rete si appoggia in modo massiccio al confronto, soprattutto con le bande più lontane dal presente, e il suo vantaggio svanisce non appena l'ordine dei passi viene alterato.

Su DVS-Lip, invece, lo stesso modulo sposta il risultato di poco più di mezzo punto, e i controlli quasi non si distinguono fra loro, non per un difetto della memoria ma perché con mezzo milione di parametri i due dataset si sono rivelati frenati da limiti opposti. La rete per THU-EACT-50-CHL non riesce nemmeno a descrivere i propri dati di addestramento, e ogni variante che la aiuta ad adattarsi migliora anche la validazione. Quella per DVS-Lip impara quasi per intero il training ma ne generalizza poco più della metà, e due modelli di accuratezza quasi identica danno la stessa risposta su poco più di metà delle clip. Anche tre addestramenti della stessa architettura, con due milioni di parametri, concordano soltanto su due clip su tre. Su THU-EACT-50-CHL la memoria relazionale fornisce alla rete la capacità temporale che le manca. Su DVS-Lip il limite è la stabilità della soluzione, che nessuna memoria del secondo patch embedding cambia.

Su DVS-Lip, come la diagnosi suggeriva, la prima leva efficace è stata la distillazione, che raccoglie in un solo modello compatto ciò che varianti diverse sanno separatamente. Porta l'IRM da $57{,}5$ a $60{,}9$ di Macro-F1 in validazione e ne alza di quasi tre punti e mezzo l'accuratezza sul test, senza alcun costo in inferenza, e rende più simili fra loro le soluzioni di addestramenti ripetuti. Un bersaglio soltanto più morbido, senza l'informazione di altri modelli, non basta a ottenere lo stesso effetto. La seconda leva è nata da un dubbio che la diagnosi lasciava aperto, cioè che su DVS-Lip una parte dell'informazione utile sparisse nel pooling prima di raggiungere la memoria. Lo stesso confronto antisimmetrico dell'IRM, calcolato prima del pooling fra celle vicine e con ritardi di poche decine di millisecondi, è il correlatore di Hassenstein e Reichardt quasi alla lettera, e migliora l'IRM di due punti e mezzo sul test senza alcun guadagno su THU-EACT-50-CHL. Si somma quasi per intero alla distillazione, perché agisce su una causa diversa, e insieme le due leve portano il modello compatto dal $50{,}8$ al $56{,}8\%$ sul test. Su THU-EACT-50-CHL, dove la diagnosi indica un limite di adattamento e non di varianza, la configurazione resta l'IRM così com'è. In un modello compatto architettura e addestramento vanno scelti insieme al compito.

Oltre il regime compatto entrambi i dataset migliorano con la capacità, e su THU-EACT-50-CHL il margine dell'IRM sul riferimento della stessa larghezza scende da quasi dieci punti con mezzo milione di parametri a tre con un milione e a due con due milioni, perché con abbastanza capacità anche il riferimento riesce ad adattarsi ai dati. La memoria relazionale conta soprattutto dove il budget è scarso, che è proprio il regime per cui è stata pensata. Su DVS-Lip capacità, relazioni locali e distillazione si sommano, e il modello migliore, con due milioni di parametri, arriva al $60{,}2\%$ sul test e al $67{,}8\%$ di Macro-F1 in validazione, i valori più alti ottenuti su questo dataset. La distillazione aiuta allo stesso modo anche il riferimento, ma a parità di ricetta le relazioni locali restano avanti di oltre due punti, mentre l'IRM da solo non se ne distingue più. Il costo della memoria, infine, si lascia ridurre dopo l'addestramento, poiché le quattro matrici con cui le bande vengono riportate nei canali si esprimono in una base comune di rango ridotto, scelto con una regola fissata in anticipo, senza riaddestrare la rete e senza perdite misurabili sul test dei modelli finali. Così compresso, l'IRM costa meno del mixer a ritardi che sostituisce a ogni larghezza su THU-EACT-50-CHL, fino a circa un quinto del suo costo con due milioni di parametri, e su DVS-Lip con uno e due milioni, mentre il modello compatto distillato ne resta di poco sopra. È, in piccolo, la stessa lezione da cui il lavoro era partito. Il ramo locale, invece, si paga, e la configurazione finale di DVS-Lip esegue fra l'8 e il $23\%$ di moltiplicazioni in più del modello di riferimento, in cambio di cinque o sette punti sul test.

Su THU-EACT-50-CHL questa catena di scelte porta oltre lo stato dell'arte, e già con $0{,}46$ milioni di parametri la rete con l'IRM compresso supera di quasi diciannove punti la migliore rete spiking pubblicata. Con uno e con due milioni supera tutti i metodi che non usano dati esterni, fino al $74{,}3\%$ di accuratezza media sul test, contro il $69{,}85\%$ del migliore di essi, che usa sei volte più parametri. Soltanto un modello pre-addestrato su un dataset esterno molto più ampio va ancora oltre. Il risultato appartiene all'intera catena, dalla rete riallocata al clock, alla capacità e alla memoria, ma l'IRM vi è decisivo proprio a basso budget, dove porta la rete compatta dal $56$ al $66\%$ e, con un milione di parametri, oltre i metodi pubblicati che il riferimento della stessa larghezza non raggiunge. Su DVS-Lip il modello finale arriva allo stesso $60{,}2\%$ di Spiking MSTP, una delle due reti spiking pubblicate su questo dataset, ma i valori assoluti restano lontani dai migliori lavori, e il dataset ha avuto soprattutto il ruolo di un caso diagnostico, che ha mostrato quando una memoria a lungo termine non basta e quali leve funzionano al suo posto.

I confini di questi risultati vanno dichiarati con chiarezza, a cominciare dal fatto che quasi tutti i confronti sull'architettura si basano su un solo addestramento per configurazione, e soltanto con due milioni di parametri l'addestramento è stato ripetuto. Il confronto appaiato tiene conto di quali campioni la validazione contenga, ma non della variabilità tra un addestramento e l'altro, che su DVS-Lip è ampia, e i valori di $p$ non sono corretti per i molti confronti eseguiti. Le differenze nette osservate su THU-EACT-50-CHL reggono a questa cautela, quelle di un punto su DVS-Lip no. Modelli, varianti e ranghi di compressione sono stati scelti sulla validazione di sviluppo. Il test ufficiale, aperto soltanto a posteriori, prima per diagnosticare modelli già scelti e poi per confrontare quelli finali con la letteratura, offre una conferma importante, ma non la misura su dati mai osservati che darebbe un protocollo più rigido. I modelli, inoltre, sono addestrati su una parte del training ufficiale, e le cifre della letteratura sono quelle riportate dai rispettivi lavori, non riprodotte. Il beneficio della memoria relazionale dipende dal compito e dalla scala, netto dove la rete compatta fatica ad adattarsi ai dati e il passato offre qualcosa da confrontare, piccolo su DVS-Lip nel secondo patch embedding, dove serve invece il confronto locale prima del pooling, assente su DVS128 Gesture e CIFAR10-DVS nelle prove svolte. Il ramo locale, a sua volta, è stato provato su THU-EACT-50-CHL soltanto nel modello compatto. Anche le alternative scartate lo sono soltanto nelle versioni provate, e nulla esclude che una memoria più lunga o un orologio adattivo, costruiti in un altro modo, possano funzionare. Quanto al costo, i valori riportati sono conteggi analitici di parametri, operazioni e stato, accompagnati da una stima aritmetica dell'energia. Indicano dove il costo si colloca, non quanto varrà una volta che il modello sarà realizzato.

Nessuna delle leve esaminate, poi, ha reso la rete invariante al ritmo, e un movimento eseguito al doppio o alla metà della velocità la porta ancora fuori da ciò che conosce. Su DVS-Lip, allo stesso modo, gli scambi fra le due parole di una coppia simile, come \emph{spend} e \emph{spent}, non diminuiscono con la memoria, le relazioni locali, la capacità o la distillazione, che riducono invece tutti gli altri errori. Le clip più brevi, infine, restano le più difficili, anche se in parte per le parole che contengono più che per la loro durata.

La realizzazione su hardware è la direzione più naturale per proseguire, ed è quella a cui il capitolo~\ref{chap:fondamenti_stato_arte} rimandava. Il modello compatto si presta per alcuni aspetti a un'implementazione su FPGA, dato che la memoria relazionale opera nel secondo patch embedding, dove le posizioni sono poche, e il suo stato è un buffer circolare degli ultimi passi, con ritardi fissi e quindi indirizzi noti in anticipo. Le medie di banda si aggiornano con poche somme e sottrazioni per passo, il confronto aggiunge una manciata di prodotti elemento per elemento, e la compressione porta il buffer al livello di quello del mixer a ritardi originale o al di sotto. Il ramo locale di DVS-Lip è la parte meno favorevole, perché conserva le proiezioni degli ultimi passi su tutta la mappa a piena risoluzione e aggiunge un quinto delle moltiplicazioni nel modello compatto. Restano però difficoltà che i conteggi non vedono, perché la memoria lavora su correnti a più bit e non su spike, e richiede quindi moltiplicazioni vere, e lo stato persistente va letto e riscritto a ogni passo, con accessi in memoria che possono pesare più dell'aritmetica. La rappresentazione a lunghezza variabile, infine, rende la latenza proporzionale alla durata della registrazione. Una sintesi effettiva, con la misura di risorse, latenza e potenza, e uno studio della precisione numerica con cui rappresentare pesi, correnti e buffer sono i passi che separano questi risultati da un sistema reale.

Anche sul piano del modello restano domande aperte, a cominciare dal costo del ramo locale, che oggi lavora sull'intera mappa a piena risoluzione e potrebbe forse conservare meno celle o meno componenti senza perdere il suo guadagno. Resta da vedere, inoltre, se il ramo aiuti THU-EACT-50-CHL con una rete più grande, dove il modello è stato provato soltanto in forma compatta. I modelli distillati, poi, si lasciano comprimere meno degli altri, e resta da capire se convenga distillare in un modello già fattorizzato o comprimere prima di distillare. L'assenza di invarianza al ritmo suggerisce di esporre la rete, durante l'addestramento, a velocità molto diverse, oppure di darle una memoria capace di leggere il passato su più scale insieme. Le scale stesse, oggi ereditate dai ritardi del mixer e riscalate per ciascun dataset, meriterebbero infine un protocollo che ne verifichi la trasferibilità da un compito all'altro.

Il percorso lascia un'indicazione che va oltre i singoli numeri, cioè che in uno Spiking Transformer compatto la capacità temporale non si ottiene allargando la rete né allungando la memoria, ma decidendo con cura che cosa sia un passo, quale passato conservare e come metterlo a confronto con il presente. Lo stesso confronto orientato, ispirato al correlatore di Hassenstein e Reichardt, si è rivelato utile in entrambi i compiti, ma alla scala che ciascuno richiede, su distanze di secondi per le azioni e fra celle vicine, su poche decine di millisecondi, per le labbra. Fatte queste scelte, una memoria piccola e locale basta a cambiare il risultato, e può costare meno di quella che sostituisce. Serve però anche saper riconoscere quando il limite non è la memoria ma la stabilità della soluzione, perché in quel caso la leva giusta sta nell'addestramento più che nell'architettura.
